
\documentclass[11pt,aspectratio=169]{beamer}

%---------------------------------------
% Presentation style
%---------------------------------------
\usetheme{Boadilla}
\usecolortheme{default}

% Remove navigation buttons
\setbeamertemplate{navigation symbols}{}

% Show slide numbers
\setbeamertemplate{footline}[frame number]

%---------------------------------------
% Packages
%---------------------------------------
\usepackage{amsmath}
\usepackage{graphicx}

%---------------------------------------
% Title information
%---------------------------------------
\title{My First Beamer Presentation}
\subtitle{An Introduction to the Topic}
\author{Student Name}
\institute{University Name}
\date{\today}

%---------------------------------------
% Begin presentation
%---------------------------------------
\begin{document}
	
	% Slide 1: Title
	\begin{frame}
		\titlepage
	\end{frame}
	
	% Slide 2: Outline
	\begin{frame}{Outline}
		\tableofcontents
	\end{frame}
	
	%---------------------------------------
	\section{Introduction}
	%---------------------------------------
	
	% Slide 3: Bullet points
	\begin{frame}{Introduction}
		
		\begin{itemize}
			\item This is the first point.
			\item This is the second point.
			\item This is the third point.
		\end{itemize}
		
	\end{frame}
	
	%---------------------------------------
	\section{Main Discussion}
	%---------------------------------------
	
	% Slide 4: Text and blocks
	\begin{frame}{Main Idea}
		
		Here is a short explanation of the topic.
		
		\begin{block}{Important Concept}
			This is an example of a highlighted
			concept or definition.
		\end{block}
		
	\end{frame}
	
	% Slide 5: Equations
	\begin{frame}{An Example Equation}
		
		We can write a simple equation:
		
		\begin{equation}
			Y = \alpha + \beta X + \varepsilon
		\end{equation}
		
		where:
		\begin{itemize}
			\item $Y$ is the dependent variable.
			\item $X$ is the independent variable.
			\item $\varepsilon$ is the error term.
		\end{itemize}
		
	\end{frame}
	
	% Slide 6: Two columns
	\begin{frame}{Comparing Two Ideas}
		
		\begin{columns}
			
			\begin{column}{0.48\textwidth}
				\textbf{Idea One}
				\begin{itemize}
					\item First point
					\item Second point
				\end{itemize}
			\end{column}
			
			\begin{column}{0.48\textwidth}
				\textbf{Idea Two}
				\begin{itemize}
					\item First point
					\item Second point
				\end{itemize}
			\end{column}
			
		\end{columns}
		
	\end{frame}
	
	%---------------------------------------
	\section{Conclusion}
	%---------------------------------------
	
	% Slide 7: Conclusion
	\begin{frame}{Conclusion}
		
		\begin{enumerate}
			\item First takeaway
			\item Second takeaway
			\item Third takeaway
		\end{enumerate}
		
	\end{frame}
	
\end{document}

